% ============================================================
%  Relatório Técnico — Centro de Distribuição Automatizado
%  Formato SBC (Sociedade Brasileira de Computação)
%
%  INSTRUÇÕES:
%  1. Baixe o template oficial em https://www.sbc.org.br/documentos-da-sbc/summary/169-templates-para-artigos-e-capitulos-de-livros/878-modelosparapublicaodeartigos
%     e coloque sbc-template.sty nesta pasta.
%  2. Compile: pdflatex relatorio.tex
% ============================================================

\documentclass[12pt]{article}

\usepackage{sbc-template}       % baixar da SBC
\usepackage[brazil]{babel}
\usepackage[utf8]{inputenc}
\usepackage[T1]{fontenc}
\usepackage{graphicx}
\usepackage{url}
\usepackage{booktabs}
\usepackage{listings}
\usepackage{xcolor}

\lstset{
  basicstyle=\ttfamily\small,
  breaklines=true,
  frame=single,
  backgroundcolor=\color{gray!10},
}

\sloppy

% ── Metadados ────────────────────────────────────────────────
\title{Centro de Distribuição Automatizado:\\
       Simulação Concorrente com Raylib}

\author{Helenna Massarra Paes}

\address{Instituto Brasileiro de Ensino, Desenvolvimento e Pesquisa (IDP)\\
  Ciência da Computação e Engenharia de Software\\
  Disciplina: Programação Concorrente --- 2026/1
  \email{helenna.massarra@idp.edu.br}
}

% ── Documento ────────────────────────────────────────────────
\begin{document}

\maketitle

\begin{resumo}
Este trabalho implementa a simulação de um centro de distribuição automatizado
da empresa fictícia IDP, composto por robôs coletores, robôs entregadores e uma
esteira transportadora.  A implementação em C usa \textit{pthreads} para modelar
cada entidade como uma \textit{thread} independente e \textit{mutexes} por recurso
para coordenar o acesso a células do mapa, posições da esteira e filas das estações.
A interface gráfica foi implementada com a biblioteca Raylib em substituição ao
ncurses, habilitando o bônus previsto no enunciado.  Três cenários estáticos e uma
configuração dinâmica validam o comportamento do sistema sob diferentes cargas.
\end{resumo}

% ─────────────────────────────────────────────────────────────
\section{Organização das \textit{Threads}}
\label{sec:threads}

O programa cria $N + M + 3$ \textit{threads} ao total, onde $N$ é o número de
robôs coletores e $M$ o de entregadores.  A Tabela~\ref{tab:threads} resume as
\textit{threads} e seus argumentos.

\begin{table}[h]
\centering
\caption{Threads do sistema}
\label{tab:threads}
\begin{tabular}{llcc}
\toprule
\textbf{Thread}       & \textbf{Função}       & \textbf{Qtd.} & \textbf{Argumento}         \\
\midrule
Gerador de pacotes    & \texttt{thread\_gerador}    & 1 & \texttt{Simulacao *}       \\
Esteira               & \texttt{thread\_esteira}    & 1 & \texttt{Simulacao *}       \\
Robô coletor          & \texttt{thread\_coletor}    & N & \texttt{ArgRobo *}         \\
Robô entregador       & \texttt{thread\_entregador} & M & \texttt{ArgRobo *}         \\
Interface Raylib      & \texttt{thread\_interface}  & 1 & \texttt{Simulacao *}       \\
\bottomrule
\end{tabular}
\end{table}

Todas são criadas em \texttt{main.c} com \texttt{pthread\_create} e aguardadas
com \texttt{pthread\_join}.  A janela Raylib é aberta \emph{dentro} de
\texttt{thread\_interface} via \texttt{InitWindow}, pois a biblioteca exige que
o contexto OpenGL seja criado e usado pela mesma \textit{thread}.

O encerramento é cooperativo: quando a flag \texttt{sim.rodando} vai a zero,
todos os laços de \textit{thread} terminam naturalmente e os
\texttt{pthread\_join} em \texttt{main.c} retornam.

% ─────────────────────────────────────────────────────────────
\section{Recursos Compartilhados}
\label{sec:recursos}

A Tabela~\ref{tab:recursos} lista os recursos compartilhados e seus acessores.

\begin{table}[h]
\centering
\caption{Recursos compartilhados do sistema}
\label{tab:recursos}
\begin{tabular}{p{3.5cm}p{2.5cm}p{5.5cm}}
\toprule
\textbf{Recurso}                         & \textbf{Tipo}          & \textbf{Acessado por}                                                        \\
\midrule
\texttt{Mapa.celulas[].ocupante}         & grade de células       & todos os robôs (r/w de posição) e interface (leitura)                        \\
\texttt{Esteira.posicoes[]}              & buffer linear          & coletores (inserção), esteira (avanço), entregadores (retirada), interface   \\
\texttt{Estacao.cabeca/cauda/quantidade} & fila encadeada         & gerador (enfileira) e coletores (desenfileira)                               \\
\texttt{Estatisticas}                    & contadores inteiros    & gerador e entregadores (escrita); interface (leitura)                        \\
\texttt{Simulacao.rodando}               & flag de parada         & qualquer \textit{thread} pode escrever; todas leem                           \\
\bottomrule
\end{tabular}
\end{table}

% ─────────────────────────────────────────────────────────────
\section{Seções Críticas Identificadas}
\label{sec:criticas}

\textbf{Movimentação de robôs (\texttt{mapa.c}).}
Dois robôs podem tentar ocupar a mesma célula simultaneamente.
\texttt{mapa\_tenta\_mover} adquire \texttt{Mapa.mutex}, verifica se o destino
está livre, atualiza origem e destino e libera o \textit{mutex} em um único
\textit{lock}, garantindo exclusão mútua.

\textbf{Esteira (\texttt{esteira.c}).}
O vetor \texttt{posicoes[]} é acessado por três \textit{threads} distintas:
\texttt{esteira\_tenta\_inserir} escreve em \texttt{posicoes[0]};
\texttt{esteira\_avancar} desloca pacotes de \texttt{[i-1]} para \texttt{[i]};
\texttt{esteira\_retirar\_saida} lê e anula \texttt{posicoes[tamanho-1]}.
Todas adquirem \texttt{Esteira.mutex} antes de qualquer acesso.

\textbf{Filas de estações (\texttt{estacao.c}).}
O gerador e múltiplos coletores disputam a mesma estação.
\texttt{estacao\_enfileirar} e \texttt{estacao\_desenfileirar} adquirem
\texttt{Estacao.mutex} para garantir que a lista encadeada e o contador
\texttt{quantidade} são modificados atomicamente.

\textbf{Estatísticas (\texttt{simulacao.c}).}
\texttt{stats\_inc\_entregues} verifica a condição de término
\emph{sob o mesmo lock}: se \texttt{pacotes\_entregues >= total\_pacotes}, chama
\texttt{simulacao\_parar} ainda dentro da seção crítica, evitando que a parada
seja acionada mais de uma vez por corrida.

% ─────────────────────────────────────────────────────────────
\section{Mecanismos de Sincronização}
\label{sec:sync}

\textbf{\texttt{pthread\_mutex\_t} --- exclusão mútua por recurso.}
Cada recurso possui seu próprio \textit{mutex} (\texttt{Mapa.mutex},
\texttt{Esteira.mutex}, \texttt{Estacao.mutex} por estação,
\texttt{Estatisticas.mutex}).  A granularidade fina permite que \textit{threads}
que acessam recursos distintos avancem em paralelo.  Nenhuma \textit{thread}
adquire dois \textit{locks} simultaneamente, eliminando o risco de \textit{deadlock}
por espera circular.

\textbf{\texttt{volatile sig\_atomic\_t rodando} --- flag de término.}
\texttt{Simulacao.rodando} é declarada \texttt{volatile sig\_atomic\_t}:
\texttt{volatile} impede que o compilador cache o valor em registrador;
\texttt{sig\_atomic\_t} garante leitura/escrita atômica em x86.  O pior caso
de corrida é uma \textit{thread} ler \texttt{1} um \textit{tick} após outra
ter escrito \texttt{0}, fazendo-a executar mais um ciclo --- aceitável para esta
simulação.

\textbf{Espera ativa com \texttt{dormir\_ms}.}
Nos pontos de espera (coletor aguardando espaço na esteira; entregador aguardando
pacote na saída), as \textit{threads} dormem \texttt{passo\_ms} e tentam novamente.
O uso de \texttt{pthread\_cond\_t} exigiria \texttt{pthread\_cond\_signal} em vários
pontos dispersos, tornando a sincronização mais difícil de auditar sem ganho prático,
pois o período de espera já é determinado pelo ritmo da simulação.

% ─────────────────────────────────────────────────────────────
\section{Decisões de Projeto}
\label{sec:projeto}

\textbf{Interface gráfica com Raylib (bônus +10\%).}
O enunciado prevê \textit{ncurses} como padrão e oferece bônus para grupos que
substituam integralmente a interface por Raylib.  Este projeto implementa a interface
com Raylib: o mapa é renderizado como grade de células coloridas e o painel lateral
exibe esteira, contadores e barra de progresso em tempo real.  A lógica concorrente
é idêntica à que seria usada com \textit{ncurses}: a \textit{thread} de interface
captura um \textit{snapshot} do estado sob \textit{mutex} e renderiza em seguida,
sem manter nenhum recurso travado durante o desenho.

\textbf{Mutex por recurso em vez de mutex global.}
Um \textit{mutex} global criaria gargalo: a \textit{thread} de interface bloquearia
todos os robôs durante a renderização.  Com \textit{mutexes} por recurso, cada robô
bloqueia apenas a célula que está tentando ocupar, e a interface adquire
\texttt{Mapa.mutex} por menos de um microssegundo (cópia de buffer).

\textbf{BFS com leitura \textit{lock-free}.}
Cada passo de cada robô recalcula o BFS a partir da posição atual, lendo
\texttt{ocupante} sem \textit{mutex}.  Dados levemente desatualizados apenas fazem o
robô tentar um passo que falha em \texttt{mapa\_tenta\_mover} (atômico); no próximo
ciclo o BFS recalcula.  Isso elimina contenção no \textit{mutex} do mapa durante o
planejamento de rotas.

\textbf{Heurística gulosa para escolha de destino.}
Coletores escolhem a estação \emph{não-vazia mais próxima} (distância de Manhattan);
entregadores escolhem o \emph{despacho mais próximo}.  Não é ótimo globalmente, mas
distribui robôs naturalmente entre estações/despachos sem comunicação entre
\textit{threads}.

\textbf{Lógica \texttt{afastar\_da\_entrada}.}
Coletores ociosos próximos à entrada da esteira poderiam bloquear o último coletor
carregado, impedindo a simulação de atingir a meta.  A função
\texttt{afastar\_da\_entrada} move o coletor ocioso para a célula vizinha mais
distante da entrada, liberando o gargalo.

% ─────────────────────────────────────────────────────────────
\section{Guia de Utilização}
\label{sec:uso}

\textbf{Pré-requisitos (Ubuntu 24.04).}
A interface gráfica exige Raylib 5.5, compilado do fonte:

\begin{lstlisting}[language=bash]
sudo apt install -y build-essential git cmake \
    libx11-dev libxrandr-dev libxi-dev \
    libxcursor-dev libxinerama-dev libgl1-mesa-dev

git clone --depth 1 --branch 5.5 \
    https://github.com/raysan5/raylib.git /tmp/raylib
cd /tmp/raylib
cmake -B build -DBUILD_SHARED_LIBS=ON -DCMAKE_BUILD_TYPE=Release
cmake --build build --target install && sudo ldconfig
\end{lstlisting}

\textbf{Compilar e executar:}

\begin{lstlisting}[language=bash]
make                # compila ./idp_sim
./idp_sim 1         # Cenario Simples    (18x12, 20 pacotes)
./idp_sim 2         # Cenario Intermediario (24x16, 40 pacotes)
./idp_sim 3         # Alta Concorrencia  (32x20, 80 pacotes)
./idp_sim 0         # Configuracao dinamica (le parametros)
\end{lstlisting}

Pressione \textbf{Q} para encerrar antecipadamente.

\textbf{Testes automatizados:}

\begin{lstlisting}[language=bash]
make teste          # valida mapa, simulacao e operacoes atomicas
make teste-bfs      # prova que o BFS desvia de obstaculo em U
make teste-stress   # ciclo completo com threads reais
\end{lstlisting}

\textbf{Via container:}

\begin{lstlisting}[language=bash]
podman build -t idp_sim .
xhost +local:
podman run --rm -it \
    -e DISPLAY=$DISPLAY \
    -v /tmp/.X11-unix:/tmp/.X11-unix \
    idp_sim 2
\end{lstlisting}

\end{document}
